% -*- coding: utf-8; mode: latex; -*-

% +++
% latex = "lualatex"
% +++

\documentclass{FITpaper}

\usepackage{graphicx}
\usepackage{flushend}

% 和文タイトル
\jtitle{空き家マッチングにおける借り手の熱意を代弁する交渉対話エージェントの試作}

% 欧文タイトル
\etitle{Prototyping a Negotiation Dialogue Agent as a Proxy for Tenant Enthusiasm in Vacant House Matching}

% 著者数
\authors{cc}

% 和文著者名（要変更）
\jauthors{%
  ○○ ○○\affmark{1} &
  ○○ ○○\affmark{2}
}

% 欧文著者名（要変更）
\eauthors{%
  Firstname Lastname &
  Firstname Lastname
}

% 所属（要変更）
\afftext{1}{名古屋工業大学大学院}
\afftext{2}{名古屋工業大学}

\begin{document}

\maketitle

%----------------------------------------------------------------------
\section{はじめに}
%----------------------------------------------------------------------

日本では空き家の増加が社会問題となっており，その活用が求められている．
総務省の住宅・土地統計調査（2023年）によると，
空き家数は約900万戸に達し過去最多を更新した．
しかし，空き家オーナー（大家）の多くは積極的に物件を貸し出そうとしているわけではなく，
使い方や借り手の人柄によっては貸してもよいと考えている層が存在する．

このような文脈で注目されるのが「逆さま不動産」という仕組みである．
従来の不動産取引では大家が借り手を選ぶのに対し，
逆さま不動産では借り手が自身の活動目的や想いを記事として公開し，
大家が「この人に貸したい」と思う相手を選ぶ（表\ref{tbl:sakasama}）．
このモデルでは，借り手の熱意がいかに大家に届くかがマッチング成立の鍵となる．

\begin{table}[htbp]
  \centering
  \caption{従来の不動産と逆さま不動産の比較}
  \label{tbl:sakasama}
  \begin{tabular}{l|l}
    従来の不動産 & 逆さま不動産 \\
    \hline
    物件情報を掲載する & 借り手の夢や想いを掲載する \\
    借り手が物件を探す & 大家が借り手に関心を持つ \\
    条件から選ぶ & 想いや使い方からつながる \\
  \end{tabular}
\end{table}

逆さま不動産では，大家が借り手の記事を読んで関心を持った場合，
借り手は大家に熱意をアピールしながら，
物件が自身の目的に適しているかを確認する交渉対話が求められる．
しかし，こうした対話を多数の大家に対して自ら行うことは
借り手にとってコストが高く，不慣れな利用者には難しい．

本研究では，借り手の記事をもとに大家との交渉対話を代行する
AIエージェントを試作する．
エージェントは借り手の熱意・目的を短く率直に大家に伝える
\textbf{感情表現}と，
借り手の目的に沿って物件の適否を確認する
\textbf{適合確認の質問}の2つの機能を担う．
特に感情表現に注力し，消極的な大家に「この人に貸したい」と
感じさせる対話の実現を目指す．

本稿の構成は以下の通りである．
2節で関連研究を述べ，3節でシナリオと課題を整理する．
4節でシステム設計，5節で実験設計を示し，
6節で結果と考察を述べ，7節でまとめる．

%----------------------------------------------------------------------
\section{関連研究}
%----------------------------------------------------------------------

\subsection{ロールプレイング言語エージェント}

大規模言語モデル（LLM）を用いてペルソナを与え，
特定の人物として対話させるロールプレイング言語エージェント（RPLA）の研究が進んでいる
\cite{chen2024rpla}．
RPLAにおける課題として，ペルソナとして与えられる情報量の少なさに起因する
事実の創作（hallucination）や，ペルソナ整合性の低下が挙げられている．

本研究が対象とする借り手記事は，
特定個人の目的・価値観を記述した
\textbf{Individualized Persona}~\cite{tseng2024twotales}
に近いが，情報量が限られるため，
記事にない事実を発話するリスクへの対処が必要である．
本研究では制約リスト（constraints）を構造化プロファイルに明示することで
この問題に対処する．

\subsection{目的志向対話の評価}

多ターン対話エージェントの評価では，
ペルソナ整合性（Persona Consistency），
根拠忠実性（Faithfulness），
タスク達成（Task Success）など複数の軸が用いられる
\cite{chen2024rpla}．
本研究ではこれらに加え，
借り手の熱意が大家に届いているかを総合的に測る
自然さ（Naturalness）を評価軸として設定する．

\subsection{不動産・マッチング領域へのAI応用}

不動産領域へのAI応用としては，ユーザの希望条件に基づく物件推薦や
問い合わせ対応チャットボットの研究が存在する．
これらは大家側から借り手へ情報を提供する方向の設計が中心であり，
借り手が大家に対してアピールする方向の対話を設計した研究は少ない．
逆さま不動産のような借り手主導のマッチング場面において，
借り手の熱意を代弁する交渉対話エージェントを試作した研究は，
著者の知る限り先行事例がない．

%----------------------------------------------------------------------
\section{対象シナリオと課題}
%----------------------------------------------------------------------

\subsection{逆さま不動産における借り手の立場}

逆さま不動産では，借り手は記事として自身の夢・目的・背景を公開している．
大家がその記事を読み，興味を持った場合に交渉が始まる．
交渉においては，借り手が熱意を積極的にアピールしながら，
物件が自身の目的に適しているかを確認するやり取りが求められる．

しかし，こうした交渉的な対話を多数の大家に対して自ら行うことは
借り手にとってコストが高く，慣れない利用者には難しい．
そこで，借り手の記事やプロファイルをもとに大家との対話を代行する
AIエージェントが有用と考える．

\subsection{システムの構図}

本研究が想定するシステムでは，
1つの物件に対して複数の借り手エージェントがそれぞれ独立に
熱意をアピールしながら対話する（1対多の構図，図\ref{fig:system}）．
大家は「この人に貸したい」と感じた借り手と対面での詳細交渉に進む
（本研究のスコープ外）．

\begin{figure}[htbp]
  \centering
  % TODO: システム図を作成して挿入
  % \includegraphics[width=\columnwidth]{system_figure}
  \fbox{\rule{0pt}{3cm}\rule{\columnwidth}{0pt}}
  \caption{1対多の対話構図（大家1人に対し複数の借り手エージェントが対話）}
  \label{fig:system}
\end{figure}

\subsection{借り手エージェントに求める機能}

借り手エージェントは次の2つの機能を担う．

\textbf{感情表現（主要機能）：}
借り手の夢・目的・想いを短く率直に大家に届け，
「この人に貸したい」と感じさせる発話を生成する．
長文の経歴紹介ではなく，核心となる動機を最初の発話で伝えることを重視する．

\textbf{適合確認の質問（補助機能）：}
「○○に使いたいのですが，○○はありますか？」のように，
借り手の目的に照らして物件が適しているかを確認する質問を自発的に行う．

%----------------------------------------------------------------------
\section{提案システム}
%----------------------------------------------------------------------

\subsection{システム概要}

本システムは，借り手AI（2条件）と大家LLMが交互に発話する
多ターン対話実験基盤である．
対話は最大6往復とし，ログをJSONで保存する．
モデルにはClaude Haiku（claude-haiku-4-5，temperature=0）を使用する．

\subsection{借り手エージェントの設計}

借り手エージェントへの入力形式の違いによって，
以下の2条件を設定し比較する（表\ref{tbl:conditions}）．

\begin{table}[htbp]
  \centering
  \caption{比較条件}
  \label{tbl:conditions}
  \begin{tabular}{l|p{5cm}}
    条件 & 入力・特徴 \\
    \hline
    Baseline & 借り手記事の全文をそのままプロンプトに入力 \\
    Structured & 目的・希望・制約・質問リストに整理したプロファイルを入力．感情表現の指示を明示的に追加 \\
  \end{tabular}
\end{table}

Structured条件では，初発話で「なぜこの場所でやりたいのか」
という動機を必ず伝えるよう指示し，
大家が消極的な反応をしても諦めずに想いを伝え続けることを
プロンプトで明示する．
また，1回の発話を3〜4文以内に制限し，
読まれやすい短い発話を促す．

表\ref{tbl:dialogue_ex}に，
陶芸家の借り手ケースにおける両条件の初発話の例を示す．
Baseline条件では記事にある情報を並列的に述べるのに対し，
Structured条件では動機を先に示してから質問へつながる構造になっている．

\begin{table}[htbp]
  \centering
  \caption{初発話の例（陶芸家ケース，大家の導入発話後）}
  \label{tbl:dialogue_ex}
  \begin{tabular}{l|p{5.2cm}}
    条件 & 借り手の初発話 \\
    \hline
    Baseline & 静岡で陶芸をしている藤井と申します。住居兼アトリエとして一軒家を探しています。120\,m$^2$程度で，ガス窯を置けるスペースと駐車場があればと考えております。 \\
    \hline
    Structured & 海外からも注文が届くようになった今，落ち着いて窯と向き合える場所を求めてまいりました。ガス窯で焼き締めの陶器を制作しており，煙突を設置できる離れや倉庫はございますか？\\
  \end{tabular}
\end{table}

\subsection{大家LLMの設計}

大家LLMは，物件情報（property\_facts）のみを根拠に回答する
grounding 型のLLMである．
大家のスタンスは「積極的に貸したいわけではないが，
人柄や使い方によっては貸してもよい」として設定し，
過度な歓迎や売り込みを抑制する．

また，物件情報の中に代替手段がある場合は
「○○はないが，代わりに○○がある」のような
代替案の提示も行う．
例えば，都市ガス未接続の場合でも
「都市ガスは未引き込みですが，プロパンガスは利用可能です」
と回答する設計とした．
これにより，借り手は物件の可能性を柔軟に判断できる．

%----------------------------------------------------------------------
\section{実験}
%----------------------------------------------------------------------

\subsection{実験設定}

\textbf{使用モデル：}
借り手LLM・大家LLMともに claude-haiku-4-5
（temperature=0，最大6往復）を使用した．

\textbf{物件：}
三重県桑名市に所在する架空の古民家を設定した
（木造2階建て・延床面積約115\,m$^2$，敷地内に離れ約25\,m$^2$，家賃月額4.8万円）．
物件情報（property\_facts）として，広さ・ガス・電気・水道・DIY可否など
16項目を設定した（表\ref{tbl:property}）．
代替案提示の機会が生まれるよう，
都市ガス未引き込みだがプロパン利用可能，
100Vのみだが200V工事は要相談（借り手負担），
など制約と代替が対になる項目を意図的に設計した．

\begin{table}[htbp]
  \centering
  \caption{物件情報（property\_facts）の主要項目}
  \label{tbl:property}
  \begin{tabular}{l|p{4.8cm}}
    項目 & 内容 \\
    \hline
    離れ & 約25\,m$^2$。アトリエとして利用可能（煙突工事は要確認） \\
    ガス & 都市ガス未引き込み，プロパン利用可 \\
    電気 & 100V。200V工事は要相談（借り手負担） \\
    水道 & あり。キッチン・トイレ・風呂完備 \\
    DIY & 可（退去時原状回復前提） \\
    家賃 & 月額4.8万円（管理費込み） \\
    駐車場 & 敷地内2台分 \\
  \end{tabular}
\end{table}

\textbf{借り手ケース：}
逆さま不動産の掲載記事から1ケース（陶芸家）を選出しパイロット実験を行った．
記事はそのままBaseline条件の入力とし，
Structured条件では目的・希望（5項目）・制約（3項目）・
質問リスト（5項目）をLLM補助＋研究者確認で構造化した．
Task Successの評価のために，借り手が大家から引き出すべき
5項目のスロット（表\ref{tbl:slots}）を設定した．

\begin{table}[htbp]
  \centering
  \caption{評価用スロット（借り手が引き出すべき情報）}
  \label{tbl:slots}
  \begin{tabular}{l|p{5cm}}
    スロットID & 内容 \\
    \hline
    kiln\_space & ガス窯スペースの有無・煙突設置可否 \\
    gas\_supply & ガスの引き込み状況（都市ガス/プロパン） \\
    area & 住居＋アトリエの延床面積 \\
    rent & 家賃と契約条件 \\
    parking & 駐車場の台数 \\
  \end{tabular}
\end{table}

\subsection{評価軸と評価方法}

以下の4軸で人手評価を行う（表\ref{tbl:rubric}）．
評価者は研究者本人1名とし，対話ログをWebビューアで閲覧しながら採点した．

\begin{table}[htbp]
  \centering
  \caption{評価軸}
  \label{tbl:rubric}
  \begin{tabular}{l|p{4.5cm}}
    評価軸 & 定義 \\
    \hline
    Faithfulness & 記事にない事実を発話していないか（違反回数） \\
    Persona Consistency & 借り手の目的・価値観と合っているか（5段階） \\
    Task Success & 必要な物件情報（スロット）を聞き出せたか（達成スロット数） \\
    Naturalness & 大家との会話として自然か（5段階） \\
  \end{tabular}
\end{table}

Task Successの評価は，各スロットについて
大家が対応する回答を返したターンを根拠として
人手でチェックリストに記録した．
借り手が質問したのみで大家が未回答の場合はカウントしない．

%----------------------------------------------------------------------
\section{結果と考察}
%----------------------------------------------------------------------

% TODO: 実験実施後に記入

\subsection{条件間の比較}

% TODO: Faithfulness・Persona Consistency・Task Success・Naturalness
% 各評価軸のスコアを条件ごとに表にして記入

\subsection{感情表現の分析}

% TODO: Structured 条件で感情表現が有効に機能した例を対話ログから抜粋
% 例：大家の消極的な返答に対して借り手が諦めずに想いを伝え続けたターンを紹介

\subsection{大家LLMの応答分析}

% TODO: 代替案提示が発生したターン数・内容を記入
% 例：「都市ガスはないがプロパンは利用可能」など

\subsection{考察}

% TODO: Structured 条件の有効性と限界
% 感情表現が対話品質（Naturalness・Persona Consistency）に与えた影響
% 代替案の提示が Task Success に与えた影響

%----------------------------------------------------------------------
\section{おわりに}
%----------------------------------------------------------------------

本研究では，逆さま不動産における借り手の熱意を代弁する
交渉対話エージェントを試作した．
借り手記事をそのまま入力するBaseline条件と，
感情表現の指示を含む構造化プロファイルを入力するStructured条件を比較し，
入力形式の違いが対話品質に与える効果を分析した．

Structured条件では，動機を先に伝えてから確認質問へつなぐ構造的な発話が
初発話から一貫して現れ，Baseline条件とは質的に異なる感情表現が得られた．
また，大家LLMに代替案の提示機能を組み込むことで，
借り手が物件の可能性を柔軟に判断できる対話環境が実現した．

今後の課題として，評価の拡充（研究室メンバーによるアンケート評価），
構造化プロファイルの自動生成，
より多くの借り手ケースでの実験が挙げられる．

\acknowledgment{%
% TODO: 謝辞
}

\begin{thebibliography}{9}

\bibitem{chen2024rpla}
  Y.~Chen, et al.,
  ``From Persona to Personalization: A Survey on Role-Playing Language Agents,''
  arXiv:2404.18231, 2024.

\bibitem{tseng2024twotales}
  B.-H.~Tseng, et al.,
  ``Two Tales of Persona in LLMs: A Survey of Role-Playing and Personalization,''
  arXiv:2406.01171, 2024.

\end{thebibliography}

\end{document}
